
\def\PUDcodigo{EME102}

\def\PUDcodacad{04506.53}

\def\PUDdisciplina{Tribologia e Lubrificação}

\def\PUDsemestre{Opt}

\def\PUDhoras{80}

%NOVOS ---------------------------------------------------
\def\PUDhorasTeoria{60}
\def\PUDhorasPratica{20}

\def\PUDinterECA{
\hyperref[PUD:EME107]{Metrologia (04506.14)}

\hyperref[PUD:EME117]{Processos de Fabricação II (04506.30)}

\hyperref[PUD:EME110]{Materiais II (04506.23)}

\hyperref[PUD:EME116]{Manutenção Industrial (04506.33)} 

\hyperref[PUD:EME125]{Análise de Falhas (04506.44)}
}

\def\PUDatuaisECA{---}

\def\PUDrecursos{Projetor multimídia; Quadro branco e pincel.}

%----------------------------------------------------------

\def\PUDcreditos{4}

\def\PUDprerequisito{ \hyperref[PUD:EME108]{EME108} }

\def\PUDpreacad{ \hyperref[PUD:EME108]{Materiais I (04506.20)} }

\def\PUDobjetivos{Conhecer as propriedades das superfícies dos materiais;  Compreender os princípios que regem a interação entre superfícies em contato e movimento relativo entre si;  Compreender os princípios físicos do atrito;  Conhecer e entender os mecanismos de desgaste mecânico;  Saber aplicar o método de lubrificação e o tipo de lubrificante adequados a cada situação mecânica.}

\def\PUDementa{Topografia de Superfície;  Superfícies em Contato;  Parâmetros Superficiais;   Mecânica do Contato;  Teoria do Atrito;  Mecanismos de Desgaste: Abrasivo, Adesivo, Fadiga Superficial e Triboquímico;  Considerações de Desgaste em Projeto Mecânico;  Tipos de Lubrificantes;  Análise de Óleos Lubrificantes.}

\def\PUDconteudo{ & Propriedades dos materiais de engenharia. 
\\ 
 & Topografia de superfície. 
\\ 
 & Introdução à mecânica do Contato. 
\\ 
 & Teorias do atrito. 
\\ 
 & Introdução ao desgaste dos materiais.  
\\ 
 & Introdução aos lubrificantes industriais. }

\def\PUDmetodo{Aulas expositórias que podem ser teóricas e/ou práticas, onde as práticas no laboratório serão marcadas no decorrer da disciplina.}

\def\PUDavalia{A avaliação será feita com aplicação de provas teóricas e/ou práticas, além da possibilidade de inclusão de trabalhos e seminários no decorrer da disciplina.}

\def\PUDbibbasica{ & \href{www.biblioteca.ifce.edu.br/index.asp?codigo_sophia=23850}{Carreteiro}, Ronald P.  Lubrificantes e Lubrificação Industrial.  1ª edição.  Rio de Janeiro - RJ.  Editora INTERCIÊNCIA.  2006.  503p.  ISBN: 8571931585. 
\\
 &  \href{www.biblioteca.ifce.edu.br/index.asp?codigo_sophia=67018}{Arnell}, R. D.  Tribology: principles and design applications.  London: Macmillan Education, 1991.  255p.  ISBN: 9780333458686.
 \\
 & \href{www.biblioteca.ifce.edu.br/index.asp?codigo_sophia=23817}{Duarte Júnior}, Durval.  Tribologia, Lubrificação e Mancais de Deslizamento.  1ª edição.  Rio de Janeiro - RJ.  Editora CIÊNCIA MODERNA.  2005.  239p.  ISBN: 8573933283. 
}

\def\PUDbibcomplementar { &  \href{www.biblioteca.ifce.edu.br/index.asp?codigo_sophia=62087}{Affonso}, Luiz Otávio Amaral.  Equipamentos mecânicos: análise de falhas e solução de problemas.  3º ed.  Rio de Janeiro, RJ: Qualitymark, 2014.  387p.  ISBN: 9788541400367.
 \\
 & \href{www.biblioteca.ifce.edu.br/index.asp?codigo_sophia=23804}{Silva}, Napoleão Fernandes da.  Compressores alternativos industriais: teoria e prática.  Rio de Janeiro, RJ : Interciência, 2009.  ISBN: 9788571932159. 
\\ 
& \href{www.biblioteca.ifce.edu.br/index.asp?codigo_sophia=23969}{Nepomuceno}, L. X.  Técnicas de Manutenção Preditiva – Volume 1.  1ª edição.  São Paulo - SP.  Editora EDGARD BLUCHER.  1989.  501p.  ISBN: 9788521200925. 
\\
 &  \href{www.biblioteca.ifce.edu.br/index.asp?codigo_sophia=40282}{Norton}, Robert L.  Projeto de máquinas: uma abordagem integrada.  4º ed.  Porto Alegre, RS : Bookman, 2013.  ISBN 9788582600221. 
\\ 
 &  \href{biblioteca.ifce.edu.br/index.asp?codigo_sophia=64982}{Ashby}, Michael. Seleção de materiais no projeto mecânico. Rio de Janeiro: Elsevier, 2012. 673 p., il. ISBN 9788535245219.
%  ARNELL, R.  D. , DAVIES, P. B.;  HALLING, J.;  WHOMES, T. L.  Trybology Principles and Design Applications 1ª edição Reino Unido Editora MACMILLAN ISBN 9780333458679. }
%{ &  ZUM GAHR, K.  H.;  Microstructure and Wear of Materials.  e-Book.  Editora NORTH HOLLAND.  ISBN 9780080875743. 
%\\ 
% &  HUTCHINGS, I.  M.;  Tribology, Friction and Wear of Engineering Materials, 1ª Edição.  Reino Unido.  Editora EDMUNDSBURY PRESS.  1992.  ISBN 9780340561843. 
%\\ 
 %&  ARNELL, R.  D. , DAVIES, P. B.;  HALLING, J.;  WHOMES, T. L.  Trybology Principles and Design Applications 1ª edição Reino Unido Editora MACMILLAN ISBN 9780333458679.
 }

\def\PUDautor{Venceslau Xavier de Lima Filho}

\def\PUDdata{2013-06-17}

\def\PUDrevisao{2}

\def\PUDrevisor{Venceslau Xavier de Lima Filho}

\def\PUDdatarevisao{2019-05-15}

\PUDcorpo
